% ============================================================
%  Niels Treschow, "Gives der noget Begreb eller nogen Idee om
%  enslige Ting? Besvaret med Hensyn til Menneskeværd og Menneskevel."
%  Det Kongelige Danske Videnskabers Selskabs Skrifter, femte Deel (V),
%  1. Hæfte. Signed "1807"; volume published Kjøbenhavn 1810.
%  TRANSLATION (English) — Hans Halvorson
%
%  Source of truth: transcription.tex (Danish, COMPLETE & image-verified).
%  Translate FROM it, not from the scan. Printed pp. 225--254 = PDF 235--264.
%
%  KEY TERMINOLOGY (fixed; keep consistent — see RESUME-NOTES.md):
%    enslig / enslige Ting / det Enslige  -> singular / singular things /
%        the singular          [NOT "individual" — Treschow uses Individ too]
%    Enslighed                            -> singularity
%    Individ / Individuer / Individualitet-> individual(s) / individuality
%    besynderlig                          -> particular  (18c sense, not "strange")
%    "Hielpemidler til Oversyn" (p.240)   -> "aids to survey"
%    Begreb / Idee / Forestilling         -> concept / idea / representation
%    Gienstand / Grundform / Kiendemærke  -> object / fundamental form / criterion
%    Forstand / Fornuft / Sands           -> understanding / reason / sense
%    Kierne / Kime                        -> kernel / germ
%    Slægt / Art / Afart                  -> genus / species / variety
%  Printed page numbers are carried over as "% printed p. NN" comments so the
%  translation mirrors transcription.tex 1:1.
% ============================================================
\documentclass[12pt,a4paper]{article}
\usepackage[utf8]{inputenc}
\usepackage[T1]{fontenc}
\usepackage{libertinus}
\usepackage{libertinust1math}
\usepackage{textalpha}   % for Greek glyphs, should any arise
\usepackage{enumitem}    % in case any enumerated lists are needed
\usepackage[english]{babel}
\usepackage[protrusion=true,expansion=true]{microtype}
\usepackage{setspace}
\usepackage[margin=1.2in]{geometry}
\usepackage{fancyhdr}
\usepackage{hyperref}

\hypersetup{
  pdftitle={Is There Any Concept or Any Idea of Singular Things?},
  pdfauthor={Niels Treschow},
  hidelinks
}

\pagestyle{fancy}
\fancyhf{}
\rhead{Treschow, \emph{On Singular Things} (1807/1810)}
\cfoot{\thepage}

\setstretch{1.3}

\title{\textbf{Is There Any Concept or Any Idea\\ of Singular Things?}\\[6pt]
       \large Answered with Regard to\\[4pt]
       \large Human Worth and Human Welfare}
\author{By Professor \textsc{Treschow}\\[4pt]
  \small Translated from the Danish by Hans Halvorson}
\date{\small Writings of the Royal Danish Society of Sciences,\\
      Part V, 1st Fascicle. 1807 [published Copenhagen 1810]}

\begin{document}
\maketitle
\thispagestyle{fancy}

\bigskip

% === TRANSLATION BEGINS HERE (body: printed p. 225) ===

% printed p. 225
\noindent
Before I proceed to the investigation of the question here proposed, it is
necessary to explain its meaning and likewise its importance. By concept
I understand, with many recent philosophers, more than the mere representation.
The latter may be sensuous, drawn from immediate experience, which comprehends
only the surface of things, not what is inward and abiding in them. The objects
of concepts are, on the contrary, the essential and unchangeable properties of
things; and concepts themselves are either, as genera and species, formed by
mere abstraction from experience, or have, through some other activity of the
power of thought according to its own laws, acquired a form at once fixed and
proper to them. This form consists in this, that concepts, together with a
certain matter, exhibit a relation between what is known from experience and
something unknown: likeness, for example, a property common to several things
otherwise different; cause or force, an inward constitution that contains the
ground
% printed p. 226
of what is outward and visible. Ideas, again, represent the highest principles,
which may be partly the ultimate end of our actions, partly the most general
and first real grounds of the objects themselves. Both are infinite, and
therefore attainable only by an unceasing approximation---the one in our
theoretical investigations, as we pursue the long series of causes up to its
highest member; the other in our actions, which, however excellent they may be,
still fall far short of the ideal's perfection.

Both concepts and ideas have this in common, that they contain something
abiding or essential, whereas the representations of the senses inform us only
of the changes to which we ourselves, as well as all outward things, are
subject. But singular objects are the proper aim of experience, that is, of the
senses: what is abiding and general in them the understanding or reason
appropriates to itself, and leaves to science to order, to investigate, to work
up; what remains it is wont to discard, like a useless shell whose kernel has
been picked out.

The singular human being is no exception to this rule. The sciences that make
our own species their proper aim---such as anthropology, morals, politics, and
legislation---consider only humanity in general, not the human being. The most
practical among them make the whole, the general, their end, and set the parts
aside, or hold themselves entitled to regard them as mere means to the former.
Hence so many thousands are sacrificed for a so-called general good: on no
altars does human blood smoke more often than on those consecrated to humanity
% printed p. 227
itself. For singular things, it is thought, are perishable; the species or the
genus alone does not perish. The properties by which one individual is
distinguished from another are merely empirical; but if they have no
unchangeable and essential character, one can form no fixed concept of them.
What, then, are they but a kind of mirage or phantasm, pleasant enough perhaps
to look upon for a short time? But why spare them, if they stand in the way of
any plan for eternity, seeing that they are destined to vanish in the next
moment?

Against this it has been said, on the other side, that genus and species
nevertheless do not subsist except through singular things, and that the former
must therefore perish along with the latter. This answer, which agrees entirely
with the opinion adopted by most philosophers, is perhaps less satisfying than
it is fit to stop the mouths of those who, without any deeply thought-out
theory, would defend the wrong they inflict, under the appearance of the good of
the whole, upon so many of their fellow men. To give out the species as a mere
concept without a particular object would be the easiest way of settling the
matter, if it fitted entirely all the singular appearances one means thereby to
explain. To these belongs, first, the striking likeness one cannot avoid
noticing between plants and animals of the same species; next, the forces that
in all, or in many kinds of beings, always act according to the same laws. For
the describers of nature will not admit that their natural species have any
other ground than the mind's abstraction; nor, on the supposition that natural
forces are merely general concepts, will one be able to understand how
% printed p. 228
the same force---that of attraction and repulsion, for example---can, on the
several ways of representing it that have been devised, either be found by a
kind of division in many singular things at once, or how of a force that is the
same only in respect of its species there can be as many individual ones as
there are substances possessing it. The last seemed to \emph{Aristotle}, among
several such examples, so unthinkable that, although he otherwise rejected
Plato's opinion concerning the origin of concepts and agreed with the empirical
philosophers cited above, he nevertheless assumes a single universally active
understanding, which in each human being expresses itself in a particular
manner, but is not numerically different. One only fails to grasp how he did not
see that the same must then hold just as well of every other force that lies
either in human nature or in the nature common to all.

It cannot, consequently, be regarded as so settled a matter that the concepts of
the understanding---that of humanity as a whole, for example---can have no real
objects outside the individuals themselves. The Platonic Ideas, so far as they
are supposed to be forms or patterns after which all that belong to a species
are formed like impressions, neither Aristotle, Locke, nor Kant has sufficiently
refuted---indeed, has not even expressly contradicted---although they have shown
clearly enough the incorrectness of many partly false or distorted
interpretations of this theory. In practical life, in the judgments of common
human understanding itself, in the plans of legislators and statesmen, the
influence of this opinion shows itself even when those who adopt or presuppose
it are scarcely conscious of the fact
% printed p. 229
themselves, or have formed any distinct concept of it. If one asks after the
objects of these ideas, it does not seem unreasonable either that they may lie
buried both in ourselves and in the whole of outward nature, as though beneath a
mist or a blinding glare, and that to become distinctly aware of them---not
merely to divine them, as most do---a sharper sight is required, which, whether
given by nature or acquired by one's own diligence, is nevertheless found in
only a few.

But even if this theory, so far as it extends, is correct, might it not
nevertheless be far too narrow, and false, if it is moreover exclusive? Are
there not likewise both ideas and concepts of singular things, or are those of
species and genera alone possible? This question is, especially when applied to
human beings, exceedingly important. It amounts, namely, to this: has every
human being, just as much as the species itself, a certain fundamental form or
essential character that remains the same under all the circumstances of life?
Is individuality anything other than a play of changing shapes or new
determinations, in which nothing is imperishable except what each has in common
with all others, that is, with the general human nature? Or are these changes
only successive developments of a kernel which---whether it is just beginning to
grow, or has already become a fruit-bearing tree, or after several
transformations has become unrecognizable even to itself---nevertheless always
retains the original fundamental features? Is, finally, the goal at which
everyone aims, even unconsciously, the pattern after which he must be formed, in
truth an infinite idea which, notwithstanding all multiplicity,
% printed p. 230
indeed apparent conflict, has inward unity and unshakable fixity in itself? It
is clear that only in this case has every individual an indescribable worth,
whereas in the other the species alone has it. That this difference of opinions
must necessarily have especially important practical consequences it is almost
superfluous to recall.

According to the exposition of ordinary metaphysics, the criterion of
singularity consists in a complete determinateness, both in respect of every
constitution one can ascribe to things and in respect of time and place. Since
both of these are continually changing, it seems to follow that in singular
beings, as such, there is nothing unchangeable. But this criterion is also
merely logical, and says nothing more than that, in order to distinguish such
things from others of the same species, one has no other means than
descriptions, because their logical essence itself, which definitions can
determine---let alone their metaphysical and real essence---is entirely unknown
to us, whereas in the case of general concepts we at least know the first.

The Scholastics come much nearer to the main point, and their disputes over
whether the identity of singular things lies in the permanence of matter or of
form betray both acuteness and insight. The first is common to many: if the
unity of the material were enough, and form were taken into no consideration,
then the concept, whose essence is form, could not possibly remain the same. For
out of the same piece of wood there may be hewn a block, and the image of a god;
form, on the other hand, is on the whole changeable, unless there are certain
laws that determine it for every individual and restrict its changeability
within such limits that the inwardly forming power can, according to the
particular nature or
% printed p. 231
idea of singular beings, develop itself in a constant advance toward greater and
greater perfection, while these, on the contrary, in virtue of an organic
connection among all that conduces thereto, mutually awaken and help one
another.

In Plato's opinion the Ideas were, like patterns from which many copies are
taken, general concepts; and the difference of singular things rested likewise
not upon the Idea's, but upon matter's greater or lesser aptness to express it.
The more recent Platonists were of the same opinion. The Schellingian system,
which at bottom is nothing other than the Neoplatonic in a shape altered to suit
the present state of the natural sciences, likewise makes it a point of honor to
raise itself above singular things as sensuous objects, or wavering images of
the Absolute and of the Ideas, that is, of its eternal and unchangeable forms.
Individuals ought therefore themselves always to strive to return to the general
nature; and although new ones constantly step forth from it upon the stage, they
are nevertheless no longer the same. Particular souls are small parts of the
great ocean which, first transformed into vapors or mist---the emblem of the
corporeal world---and then into raindrops, gather in springs, brooks, and
rivers, in order, after a brief separated existence whose pleasures privation
and longing continually embitter, to repent of their defection and become
blessed again through forgetfulness.

If this representation is correct, then human life not only has little or no
worth, but is even a real evil, from which one must strive to become free the
sooner the better. The body, sensuousness, and all that belongs to our
individuality is a perishable garment that only hinders the mind's free
% printed p. 232
movements, a weight that presses us down to the earth, or a prison through whose
thick walls the reflected light from open nature cannot penetrate. True, we are
not permitted to break out; we must remain until we are either acquitted or
released. But how much more beneficent, on this presupposition, is the voice of
cruelty and bloodthirstiness, which, in leading thousands to the shambles, at
the same time opens to them the gates of true freedom, than that of mildness and
humanity, which, by crying forgiveness, really imprisons the sinner anew? Could
one indeed wish that this theory won general approval, and that rulers and
legislators might on that account believe themselves entitled to treat human
beings as mere means, to use them as materials for some magnificent work or
other that they intended to erect? Is not the concept of right and wrong founded
entirely upon personality, i.e.\ individuality? Abolish this; suppose it could
be understood without anything self-subsistent and abiding, that persons were
only fleeting shadows that come and pass away---how could any right then belong
to them, who in the next moment would have to be lost? How could one punish or
reward them, if they themselves vanished together with the action imputed to
them? All natural as well as civil laws presuppose something abiding in persons:
they would be unfair and unreasonable if the relation among singular beings were
not as lasting as it is determinate.

If in this matter one consults experience alone, one receives, as in most cases
where general validity is at stake, a quite ambiguous answer. In the whole of
inorganic nature there is found none of the so-called singular beings there that
% printed p. 233
possesses any unchangeable character. The species themselves are so wavering
that fixed concepts of them are very difficult to obtain. Every inorganic body
is an aggregate of many; every part can subsist by itself just as well as in the
whole without losing its proper nature. By division the form can also be
changed, which is as little essential to them as the matter: who, therefore, can
say wherein their individuality consists? When does the river or the building
cease to be the same? The example of Theseus's ship is well known and has often
been cited in application to this matter. Even were the fundamental forms of
crystals---which in this part of nature nevertheless furnish the best
criteria---unchangeable, they are still, like Anaxagoras's homoiomeries,
divisible, perhaps to infinity; hence likewise aggregates, not true singular
beings. Real individuals we find only in organic nature. The more perfect the
organization is, the greater the permanence there is in their form, and the less
can any part be dispensed with, or any exchange take place, without essential
change in the whole.

Among all figures the circle is the most perfect, because every point in it is
entirely determinate, and all of them in the same way by means of the center.
There is no difference between singular circles except in magnitude. If we think
this last unbounded---which will not do in any other figure, since in these the
lines from the center to the extremities are always of unequal magnitude, and
consequently not infinite---then of the species itself there can be no more than
a single one. All other kinds of curved and rectilinear figures can also be
different in other respects, for example in the relation of their ordinates,
sides, or angles to one another.
% printed p. 234
Not without reason, therefore, did the ancients represent
the highest real principle, the Godhead as the most perfect individuality, under
the image of a circle or a sphere.

Nature's whole striving, the general drive to perfection, has no other aim than
regularity, and with it likewise permanence in form. Moral as well as physical
perfection consists in undisturbed harmony or, in other words, complete
individuality, whereby all parts are so connected that none is less essential or
necessary than another, and that no contingency has place in it.

The drive to self-preservation is by nature directed most immediately upon the
singular, only mediately upon each being's own species. The original desire for
activity is indeed by no means merely self-interested, inasmuch as the material
it compels us to work upon and to shape according to an idea is not our own
nature alone and what belongs to it, or what serves its perfection, but that of
other human beings as well---indeed, the whole external world, provided only it
is fit to receive such formation, can be its aim. We love beauty, not because it
is found in our own person, but also when we become aware of it in any other
thing whatever; everywhere, therefore, we seek to realize its idea. Nonetheless
the ground of the pleasant impression it makes is that we thereby come to
consciousness or self-feeling, i.e.\ to the sensation of our own power, its
increase or strength.

% printed p. 235
Self-consciousness is the foremost expression of individuality, the highest
step known to us that any finite being has ascended. The more clearly and
distinctly we distinguish ourselves from others, and know what marks us off from
them, the more perfect we are. All philosophy consists in knowledge of
ourselves---not according to the common human nature alone, but especially
according to the individual one. The fruit of learning and reflection is chiefly
a kind of self-recollection. All development and resolution of concepts, all
composition and connection of them, serves only to make consciousness entirely
clear, to put us in a position to give ourselves a complete account of
ourselves. Conscience and morality rest on nothing else either. To act not from
impulse but from principles---is that not to act with distinct consciousness of
them? For even in those who are determined by feelings, fundamental principles
are at work, though according to obscure and indistinct representations. The
whole excellence of our being consists, therefore, in individuality, and the aim
of all striving is to attain it in the very highest degree. It is consequently,
in the first place, legitimate: it is neither necessary, nor useful, nor
possible to suppress it. In the second place, it is the pattern after which
everyone must form himself, the idea one ought always to seek to approach; from
which it follows that ideas can be individual, and that their objects, just as
much as the species, are eternal, infinite, and unchangeable.

One may even raise the question whether there are indeed other ideas of real
things than those adduced, and whether the general opinion---that these must
either necessarily be general concepts, or that they can at least just as well
have species
% printed p. 236
and genera for their objects as singular things---is not founded upon a
confusion of the logical and ideal essence of things with the real, which only
individuals can actually possess. The dispute as to whether genera and species
truly exist outside our thoughts is, as is well known, very old. Those who hold
them to be real appeal to the fixed boundaries that nature seems to have set
between beings of different species, boundaries that make any transition from
the one to the other impossible. As concerns the more perfect organisms in the
plant as well as the animal kingdom, the matter is held by most naturalists to
be settled; the human race especially is in their view a striking proof of it,
since even the lowest among our species is so plainly distinguished from the
most excellent among the neighboring species---or, if one will, genera---that no
confusion or intermixture seems possible here except through gross ignorance.
Nonetheless there are facts and other grounds that might awaken doubt of this. I
will not even adduce the many uncertain ones that have been employed here by
various famous men, nor the conjectures and hypotheses that have been advanced,
but not sufficiently proved, by considerable authors such as Monboddo, Hallen,
Lamarck, and even our own Fabricius. In the first place, it is nevertheless
certain that the more general concepts become, and the further they remove
themselves from the individual ones, the more their boundaries run together, and
the more difficult they become to determine. The criteria of genera are in
general more wavering than those of species, and among the former most of all
the higher ones---namely, those of orders, families, and classes, of organic as
well as inorganic bodies. As our knowledge of nature, even the merely empirical,
extends itself, we see gradually vanishing those boundaries
% printed p. 237
which we believed we had determined not arbitrarily, but according to nature's
own direction. In the second place, even when not even the longest time seems
able to displace them, it is nevertheless doubtful whether this permanence is
absolute and essential, or merely relative, since we are compelled to regard as
unshakable whatever a few thousand years---to which our experience, when it
reaches its utmost, must after all confine itself---have let stand. We find, in
the third place, varieties just as well---Negroes, for example, or among dogs
pugs and greyhounds---as proper species, which, although they have probably
become so divergent through accidental circumstances, have at last taken on an
indelible character. How, then, will one make good that the same may not hold
equally of the latter?

It is, finally, not so difficult to conceive how species may have come into
being through the general laws of attraction and affinity---which in the plant
and animal kingdoms acquire several new determinations and a higher
significance---even if from the beginning nothing but individuals were brought
forth, in many though not yet so distinctly determined gradations and forms of
the same force. For since like everywhere seeks like, and since connections,
societies, the reproductive drive, and the drive to imitation must always bring
several still nearer to one another than they already are by nature, all
singular beings must thereby at length have formed groups or clusters, which may
indeed at the boundaries stand in a double and ambiguous kinship---namely, both
to those within and to those outside the same---but of which the greatest number
must nevertheless be recognizably enough separated from those belonging to the
% printed p. 238
surrounding ones. Originally, therefore, there are only individuals; species and
genera seem, on the contrary, like societies or works of art, to be of a
different and later creation.

\emph{Linnaeus}, who by a kind of stroke of genius discovered much that in
others would have been the fruit of long reflection, held that in the beginning
there were only genera, from whose combination the species subsequently arose.
Kant's hypothesis concerning the origin of the different human races, the kernel
of which is supposed to have lain in the first human pair, has in all likelihood
had the same ground in his case, although neither of them has distinctly
developed it. I should think that both meant, or at least ought to have meant,
roughly the following by it: neither genera nor species were in the beginning as
determinate as they have since become through mixtures and combinations. The
forms were at first, as in the works of sculptors, entirely rough. Nature has,
so to speak, contented itself with drawing the great outlines; the finer ones,
as well as the final polish, have come gradually. It may be conceived, for
example, that the first animals existed in a condition in which it might have
been difficult to determine whether the form was meant to represent a human
being or an irrational animal. After that shape had begun to reveal itself, it
might still be uncertain which person in particular it was the master's intention
to portray. One does indeed find everywhere physiognomies in which the
characteristic and individual seems to be lacking. But the nearer the work came
to completion, the more recognizable the distinguishing features became. Here
too nature proceeds only gradually. But it is clear that individuality is its
% printed p. 239
utmost aim. As concerns the human race, experience teaches that culture itself
contributes to this, and that for that reason one finds among civilized peoples
a multiplicity, and likewise a determinateness, in the form of singular persons
that is never so great among barbarians and savages. It must consequently be
nature's intention to individualize itself: the particular or the individual is
not a means and instrument for the general, but the end in itself.

But even if it were settled that individuals are the real---indeed perhaps the
only---objects of ideas, has this proposition then any important consequence?
Does one thereby obtain a nearer knowledge of individuals, or of the things
themselves? Can one now form a distinct concept of it better than otherwise, and
determine it so that we are put in a position to distinguish the one from the
other? If not, what use can the result be that we have sought so laboriously?
For as easy as the empirical criteria of singular things must be to find
out---criteria sufficient for use in daily life---so difficult is the
fundamental form itself either to fix in the mind or to express in words.
Nonetheless there hovers before our eyes from it a clear though confused image:
otherwise we should not learn to know things better from acquaintance and
intuition than from mere descriptions. The art of painters and sculptors
consists in grasping these fundamental forms and in so presenting them that even
by beautification they are set in a greater light rather than obscured. Thus it
is possible, out of every human face, out of every not entirely perfect
character, to form both an ideal and a caricature without effacing the
fundamental features or making them less recognizable in the one
% printed p. 240
way than in the other. Neither the roses of youth nor the wrinkles of age alter
this form: if only we are properly attentive, we become aware of it in both. Age
and culture develop it more and more: in early childhood it is entirely
indeterminate; rawness and wildness likewise make it difficult to observe---in
which conditions individuals also seem most to resemble both those reckoned to
the same species and those belonging to different ones, human beings and apes,
for example.

It makes no difference to the main point that we are not in a position to
present this form either in descriptions or in definitions. The first are taken
from the sensations of the moment, or from the impressions things make upon our
senses, in which the essential and the accidental, the objective and the
subjective, are always mixed. For definitions, language lacks the requisite
wealth of words, since of necessity it must confine itself to designating
general concepts; and although it also has names for singular objects, it has
none for the individual constitutions that make them such. The limitation of the
human understanding must content itself with knowing things in general and in
the large. But one ought not on that account to forget that the general concepts
are only aids to survey, not to exact knowledge of things in themselves.
For in the real world there are found only individuals, and their essence is not
what is expressed by the criteria of the species or the genus, but what lies in
the eternal idea, of which innumerable modifications, showing themselves in
time, are possible. This idea is not general in the sense that it comprehends
several things; but insofar as it expresses all possible states of the same one,
one might perhaps give it this name. Properly, however, it is
% printed p. 241
right to make a distinction between this abstract representation of the singular
and the concrete one.

It is not my intention to call in doubt the fundamental principles upon which
the description of nature builds. I acknowledge, with everyone else, their
necessity and their usefulness. But one would, out of love for this science, in
my view go too far if one either held the criteria of species to be essential
and unchangeable merely for the reason that the limited experience we have does
not expressly testify against it, or if in the general consideration one wished
to conceal the uncertain, the wavering, the many exceptions and deviations that
really have place with respect to the species as well. Even were the opinion of
their absolute unchangeability completely proved, it would still not follow that
the inward character of individuals is less permanent; on the contrary, analogy,
as was remarked above, permits us to infer as follows: just as the criteria of
species are less wavering than those of genera and families, so individuals, as
the lowest concepts, are subject to still less important deviations from their
fundamental form. The proposition I here defend can therefore very well subsist
alongside the theory of the describers of nature, although this science cannot
make use of it, because the former is both impossible to apply here and, if that
could be done, would lead to a prolixity far exceeding the bounds a human
science must have. In history, on the other hand, which must necessarily concern
itself with singular persons, the proposition here treated finds both a nearer
application and a confirmation---both of them in those delineations of character
which history is not in a position to sketch without a deeply
% printed p. 242
searching glance, both into human nature everywhere and into its manifold finer
and individual forms, distinguished as these are with the stamp of eternity.

In vain, nevertheless, do I seek by these and several such grounds to convince
of the correctness of my opinion those who believe that the gulf between us and
the Absolute---which one ought after all to strive to fill---is thereby made
still wider and deeper; unless I succeed in showing that this proposition not
only can subsist alongside their identity-system, but that this system even
necessarily leads to the same result. To that end I shall now explain, as
briefly and distinctly as possible, how the relation of individuals to the One
and Absolute, as well as to the eternal ideas, seems on these philosophers' own
principles best and most easily to admit of being conceived.

In however many ways this doctrinal edifice can be, and actually has been,
presented, they all seem to me to agree in the following teachings. The highest
real principle cannot be thought otherwise than as pure and unrestricted power,
activity, or acting. But since outside itself it finds no material to act upon,
and a creation out of nothing is either self-contradictory or at any rate
altogether inconceivable, this activity cannot be compared with any other known
to us except that of the understanding or of reason, so far as the latter of
itself and \textit{a priori} brings forth concepts and ideas. Now the material
for these cannot be given from without either; the forming understanding must at
the same time be the producing power; but this it can be only through
self-intuition
% printed p. 243
---or, in other words: the material for the ideas it cannot take from
anything other than itself. Sensuousness one therefore cannot ascribe to the
very highest being, because it must always be regarded as passive and a mere
capacity for being affected. Hence these philosophers teach that there is an
intuition belonging to the understanding, of which I cannot treat in this place.
But from this it now follows that the material of all ideas is the same, and
that every difference among them rests upon their form---i.e., the many
determinations, or properly, boundaries and limitations, that the former can
receive. Every idea is nevertheless, like the circle and all other lines or
figures, infinite in a certain respect, because, namely, the magnitude is
indeterminate notwithstanding the determinate form. They are likewise eternal
and necessary: it is the very essence of things that they express. Only when
they present themselves in time and space, as the forms of the sensible world,
does their eternal nature put on, so to speak, a perishable garment. From their
heavenly dwellings they descend into these sublunary regions, where all things
are partly changeable and partly show themselves at the lowest step of their
existence, and the ideas in the most degraded shape. The indwelling power seems
either, as in the coarser matter, entirely dead, or else finds itself, as in
germinating shoots, at the first beginning of its development. This time brings
with it, which can be thought neither without changes nor without an orderly
succession that must be fixed by a rule. But this succession begins with a
point, which is the very least that is real and thinkable; it grows, and becomes
at length an infinite line, or, so far as one can represent the point as
radiating, into several, indeed innumerable ones. In the same way every force is
at the beginning weak and invisible. Life borders upon death before
% printed p. 244
it gathers strength with time. Light and warmth, the most active materials, hide
themselves, like smoldering sparks before they are kindled, under the shape of a
sluggish matter. But against the unity of form, the unchangeability of the idea,
such successive advance is by no means at variance; on the contrary, the rules
according to which it develops itself constitute a part of its essence.
Furthermore: as many different forms as are possible, so many eternal ideas must
there also be; for the understanding intuits all, its omnipotence can make them
all real, and its goodness and wisdom will surely seek to bring the ordinary
so-called real things---i.e.\ those existing in time and space, which are yet
properly only weak images of the true ones, namely the objects of the
ideas---ever nearer to these patterns.

This whole theory can have no possible application except to individuals. For in
the first place, the species is not the pattern after which these must form
themselves. Neither is it so determinate that any fixed rule can be derived from
it, nor does perfection even belong to its concept; since every singular being,
according to the general opinion, which makes the species unchangeable, must
necessarily remain within the bounds of the species, however much it may deviate
from the rule. Next, not the species alone, but the singular beings too, are
different in respect of their form. It is therefore far from the case that the
former exhaust all possible forms---which must nevertheless necessarily happen,
by virtue of the infinite understanding's richness or fruitfulness.

Without binding myself to the terminology of any author or school, I have sought
to present the principles of this doctrinal edifice
% printed p. 245
in such a way as seems to me at once the clearest and likewise the most
consistent. That many both so-called Platonists and others who profess it have,
in condemning the natural striving that has individual perfection or the
preservation of the singular being for its aim, mistaken the spirit of their own
system---the reason for this is partly that they have regarded the criteria of
individuality one and all as impermanent, without attending to the essential and
unchangeable element that contains their inward ground; and into this error they
have again been led because they have not noticed that real individuals are found
only in organic nature, but have confused with them the aggregates that make up
the inorganic. Partly they have, in order (as they thought) the better to explain
the origin of evil, regarded the going-forth of the ideas---i.e.\ their imaging or
development in time and space---as a kind of Fall, which even among the most
recent philosophers has given occasion to strange conceptions; although some
among them, indeed \emph{Plotinus} himself, have seen that an infinite and
constant activity is not possible except through finite products or creatures,
that the life of these can be nothing other than an unceasing striving after the
better, and that they must consequently all be more or less evil, i.e.\
imperfect.

Yet this sort of philosopher is perhaps less difficult to convince of this than
others. Some among them come to meet us of their own accord in this respect.
\emph{Spinoza}, whom despite his very different manner of presentation I do not
hesitate to place on the same list, assumes, for example, that every human soul
% printed p. 246
thought under the idea of eternity is imperishable---by which he has meant
nothing other than that its essence or fundamental form is through all eternities
the same, not that it merely retains the criteria of the species. \emph{Fichte},
whose doctrinal edifice, as his later writings show, comes ever nearer to that
Alexandrian and now Germanic Platonism, is likewise not of the opinion that by
laying aside one's individuality one would raise oneself to a greater perfection,
or to the perfection once unhappily forfeited. Indeed, even according to the
principles of the most recent philosophy of nature, the striving to individualize
itself is general throughout the whole of nature; and what does this word mean
but the assumption of the most determinate and unchangeable form? This drive
\emph{Schelling} seems to confuse with a desire for separation, isolation, or
independence of the Absolute, with that sinking from the world of Ideas down into
the world of the senses which is supposed to be the proper source of all physical
degeneration and moral corruption. But singularity does not consist in anything
sensuous and changeable; the fundamental form, together with the rules according
to which it develops itself in every individual, is an eternal idea. In forming
itself after its pattern the soul does not forsake its origin, but much rather
returns to it. Let me be excused for using here the turns of speech of mysticism,
which for the rest I intend neither to contest nor to defend, since it is only my
intention to show that my opinion does not conflict with this school's principles
either. That it agrees with a more empirical philosophy---indeed, that one cannot
do without it in order to grasp what experience teaches---has already been proved
above.

% printed p. 247
There is, then, a fixed concept of singular things. This concept has objective
validity or reality. It stands, therefore, quite otherwise with human souls, and
with organic bodies, which may be regarded as the outwardly intuitable images of
the former, than with rough stones or clouds and collections of vapor, of which
one can only say in general what they are, but cannot adduce any particular
distinguishing mark that may not, while they still exist, be changed into its
opposite. Meanwhile one must admit that those objects, just as much as these,
seem in the end to vanish. Individuals perish while the species endures. Let them
even in the last moment of their existence retain the same essential form as in
the first, let the idea live on---the real object is nevertheless no more. The
theory is therefore, one might conclude, if not incorrect, at least without use.
It is worth the trouble to investigate whether this is so or not.

The ideas that arise in the human understanding, the plans it sketches, have not
always any real object. Together with the drawing or the sketch the work itself
is not at once finished. The understanding may be fruitful, but the fruits rest,
so to speak, in its own lap; outside itself it seems so little able to bring
forth anything that not even will and purpose are unfailing consequences of it.
In order to realize its idea it needs, as one generally believes, not only other
but even foreign forces, which are sometimes difficult enough to move to a
cooperation harmonious with it. Even if the work has been well begun, unforeseen
obstacles may interrupt it, a hostile force break it down. Defects in the design,
which we ourselves
% printed p. 248
discover, may cause us to smash the form in pieces in order to begin anew on a
more correct plan. Most human works are houses of cards, which in the heat of the
beginning we may indeed have destined for eternal remembrance, but which we
ourselves afterwards often seek to eradicate to the last trace. But should the
same hold also of the works of nature and of providence, or is precisely the
opposite of this correct? Should there ever be lacking here either the means or
the ability to carry them out; or should many a building, laid out perhaps on an
immature plan, often be destroyed by a more enlightened understanding, in order
out of its materials to raise a more perfect one better suited to the purpose?
Whether one assumes an inward necessity or an infinite wisdom which, as concerns
the arrangement of the whole world, has laid out and carries out everything for
the best---the laws of order must be posited as unchangeable, the conflicting
forces must at length bring one another into a kind of hovering equilibrium,
confusion must resolve itself into harmony, and therein must the general tendency
consist. These propositions are indeed grounded upon, or lead at length to, an
intellectual system in the whole of nature; nevertheless experience and
observations themselves press upon every investigator so irresistible a conviction
of their correctness that even those philosophers who have not been ashamed to
reckon chance among the operative causes of nature have yet let it act in an
entirely regular manner. A great proof of this is, among others, Democritus's
innumerable worlds, in which the same individual beings constantly recur.

The fundamental forms of things are consequently, like the Platonic Ideas,
unchangeable---indeed, the only real things; and these
% printed p. 249
are, as I have shown before, not species and genera as objects of the general
concepts, at least not those alone, but individuals. That the permanence of
species has always belonged, as it were, among the common articles of faith---this
has surely had, besides experience, a higher theoretical ground as well in the
unchangeability of natural laws, without which an always incomplete induction
could yet give no certainty. This ground we now with justice transfer from the
species to the individuals. For although experience is so far from confirming it
here that the perishing of singular beings seems much rather to testify against it
every moment, this appearance really arises only from a misunderstanding, as is
clearly recognized when we remain true to the principle. It has been proved that
every individual has a fundamental form whose main outlines, under all its
alternating states, would be found indelible if our ink were deep enough. These
are, namely, visible only in the changeable material that clothes and at the same
time obscures them, just as the strokes written with the sympathetic glance
become somewhat legible when another substance has given them color; but are they
therefore less real, because of themselves they do not strike the eye, or because
this cannot happen without the form thereby suffering some change?

Most naturalists have likened the organic material out of which animals and
plants are formed in the mother's womb to a germ, or have ascribed to it a kind of
preformation, i.e.\ an inward aptitude
% printed p. 250
for the determinate form, which perhaps only after the lapse of millennia becomes
able to develop itself. The organic force is, according to the constitution of
the material, different in every body, acting not without rule but according to
rules that are for every case determined most exactly; for although outward
causes that supervene may alter its direction, they cannot so annihilate its
activity that its peculiar tendency does not become recognizable to an attentive
eye. No degeneration is so great, no monstrosity even so misshapen, that the
original form is entirely perverted. It is remarkable that both the Stoics and
other materialists have adopted this preformation system, and yet may nonetheless
have believed that at death this form is so broken or crushed that a renewal of it
can never again take place. Can it not then, after the dissolution of the coarser
body, endure just as well as it knew how, before birth, to keep itself invisibly
alive---for how long is uncertain---under so many and such great destructions? For
the greatest upheavals in nature leave the proper germs themselves unharmed,
indeed sometimes even promote their growth; the finest and often hidden
dissolutions doubtless spare them likewise, or find in these smallest bodies
perhaps a resistance corresponding to their fineness and standing in inverse
relation to their magnitude. One would much misunderstand me if one believed that
by these germs, or by the fundamental form of individuals, I understand anything
other than a disposition or a
% printed p. 251
certain determination in the forming power, which is different in every
individual. This, I believe, every investigator of nature can very well
concede---whether he adopts with Bonnet the evolution-system, or with most of the
more recent a kind of epigenesis. For the rest, it would lead me too far if I
were here to show in more detail how I represent this matter to myself. The
Stoics' \emph{rationes seminales} were likewise only capacities and dispositions
toward certain forms, by no means those forms themselves in miniature.

From the standpoint at which we have arrived by these investigations, every human
being may be regarded as one of the many possible forms that their common nature
or concept can receive. Persons are consequently, in the first place, not
insignificant small parts of a whole, but each is this whole itself under his own
peculiar form. Next, one ought not to regard individuals as many more or less
perfect impressions of the same original, of which each would have to be worth so
much the less as the multitude is greater. For when of such a multitude only a
single copy remained, little would in a certain respect be lost by the perishing
of the rest. But if, on the contrary, individuals do not constitute merely a
number in which all units are alike, then their preservation, as concerns their
essential form, cannot be less important than that of species and genera. We must
consequently suppose that no individual really perishes, but that death is in
fact only a transformation, or a preparation for one---whose beginning and
% printed p. 252
progress it is nevertheless as little permitted us to follow with our
observations as what takes place with the germ before conception.

If one regards nature merely as a collection of works of art set up for viewing,
completeness nevertheless demands that nothing be lacking. Since nature never
permits itself any leap, but the series of its products must everywhere be
unbroken, just as little can any individual be missing as any genus, species, or
variety. But nature is something more than merely such a gallery: it is a
perfectly organic whole, for which every one of the self-subsistent parts that
compose it has the same necessity or importance as the whole itself, or as all the
rest have for the whole. Even the loss of the smallest cannot be made good by any
other, because precisely this one's determinate form alone fits into the place
where it stands. One can compare with this neither the imperfect works of human
beings nor even nature's own, as these present themselves in time or for the
moment in experience. In both it does indeed happen that one part can be exchanged
for another without particular loss---indeed, that the work may thereby even be
improved; but ideas are, like the ideals of art, imperishable: and although in
striving to attain them one is constantly correcting and altering, there is
nevertheless, both in respect of the parts and of the whole, something
characteristic which furnishes a rule, for every single piece as well as for the
species, from which nature itself never dares to deviate.

% printed p. 253
\emph{Aristotle} was, as is well known, of the opinion that only the general and
the great can be the object of divine providence; that this governs the whole
according to certain laws; but that the particular, the singular things, are
either far too insignificant to be worthy of this providence, or such hindrances
to its great plan that it cannot be carried out without sacrificing them. Without
belonging exactly to the school of the Peripatetics, many others have both adopted
the same principle in theory and still more often conformed to it in their manner
of acting. In statecraft it has always been the great wheel that has determined
the machine's course. But since those who direct it regard themselves as the
representatives of the species, they find therein a ground for an exception
advantageous to themselves. This ground is nevertheless entirely false, since the
societies those men represent---as \emph{Pomponatius} has already remarked in his
book on incantations---by no means constitute the species itself, but are just as
perishable as singular persons, and, so far as they are human institutions, can
have no fixed character or true individuality.

One might believe that the opposite opinion, which I have here sought to defend,
leads in practice to a still more harmful moral egoism. This consequence would be
entirely correct if either the individual could ever perish, or if by a
disinterested striving for the perfection of others it were hindered in its own
formation. But this is not how the matter stands. The objects of ideas
% printed p. 254
are, like the ideas themselves, eternal; and from morality, which is the general
law of reason enjoining inward and outward harmony, both have all their beauty
and splendor.

Yet the application of this doctrine in our dealings with our fellows, in public
and in private life, I should perhaps have had to add at length in a popular
lecture, although it is not difficult to make. But in a learned society, where
what is easier ought to be passed over and only the most involved knots, so far as
is possible, are to be untied, it is enough to point to the application. A more
pragmatic execution, entering into the particular and the concrete, is therefore
neither necessary nor, in this place, altogether fitting.

\begin{center}\rule{0.25\linewidth}{0.4pt}\end{center}

\end{document}
